% Generated draft: requires mathematical and source review.
\documentclass[11pt]{article}
\usepackage[margin=25mm]{geometry}
\usepackage{fontspec}
\setmainfont{Latin Modern Roman}
\usepackage{amsmath,amssymb,mathtools,unicode-math}
\setmathfont{Latin Modern Math}
\usepackage{enumitem,needspace,xurl,hyperref}
\setlength{\parindent}{0pt}
\setlength{\parskip}{5pt}
\setlength{\emergencystretch}{3em}
\newcommand{\N}{\mathbb{N}}
\newcommand{\Z}{\mathbb{Z}}
\newcommand{\Q}{\mathbb{Q}}
\newcommand{\R}{\mathbb{R}}
\newcommand{\C}{\mathbb{C}}
\newcommand{\F}{\mathbb{F}}
\newcommand{\A}{\mathbb{A}}
\newcommand{\PP}{\mathbb{P}}
\newcommand{\E}{\mathbb{E}}
\newcommand{\eps}{\varepsilon}
\newcommand{\dd}{\,\mathrm{d}}
\begin{document}
\section*{UnsolvedMath ID: 30005353}
\textbf{Source problem code:} catalogue entry
\subsection{Definitions and mathematical statement}
Let \(G\) be a finite connected graph, regarded as a one-dimensional CW complex, with each edge assigned length \(1\). For a closed edge path \(c\), let \(\ell(c)\) be its number of edge traversals, counted with multiplicity. The supplied record does not specify whether its attaching cycles must be simple cycles or may be arbitrary closed edge paths. Accordingly, fix a class \(\mathcal C(G)\) of admissible attaching cycles; the two interpretations to be distinguished are all simple cycles and all nonconstant closed edge paths. The definitions and questions below are conditional on that choice.

For a finite list \(c_1,\ldots,c_m\in\mathcal C(G)\), let
\[
X(G;c_1,\ldots,c_m)=G\cup_{c_1}D^2\cup\cdots\cup_{c_m}D^2
\]
be the space obtained by attaching one \(2\)-cell along each listed path. Define
\[
\alpha_{\pi_1}^{\mathcal C}(G)=\min\left\{\sum_{i=1}^{m}\ell(c_i):m\geq 0,\ c_i\in\mathcal C(G),\ \pi_1\bigl(X(G;c_1,\ldots,c_m)\bigr)=\{1\}\right\}
\]
and
\[
\alpha_{\mathbb F_2}^{\mathcal C}(G)=\min\left\{\sum_{i=1}^{m}\ell(c_i):m\geq 0,\ c_i\in\mathcal C(G),\ H_1\bigl(X(G;c_1,\ldots,c_m);\mathbb F_2\bigr)=0\right\},
\]
where \(\mathbb F_2\) is the field with two elements and the empty list has cost \(0\). Equivalently, the homological condition says that the classes \([c_1],\ldots,[c_m]\) span \(H_1(G;\mathbb F_2)\).

The supplied question asks whether, for every real \(R>0\), there exists a finite connected graph \(G\) with \(H_1(G;\mathbb F_2)\neq 0\) such that
\[
\frac{\alpha_{\pi_1}^{\mathcal C}(G)}{\alpha_{\mathbb F_2}^{\mathcal C}(G)}>R.
\]
The restriction to graphs with nonzero first homology makes the ratio defined: for a tree both costs are \(0\).

The bounded-degree question asks whether there exists a fixed integer \(D\geq 1\) such that, for every real \(R>0\), there exists such a graph \(G\) satisfying additionally
\[
\Delta(G)\leq D,
\qquad
\frac{\alpha_{\pi_1}^{\mathcal C}(G)}{\alpha_{\mathbb F_2}^{\mathcal C}(G)}>R,
\]
where \(\Delta(G)\) denotes maximum vertex degree. Thus the degree bound must be independent of \(R\).

The notation \(\alpha_{\pi_1}^{\mathcal C}\) and \(\alpha_{\mathbb F_2}^{\mathcal C}\) makes the unresolved attaching-cycle convention explicit; identifying these quantities with the source's \(\alpha_{\pi_1}\) and \(\alpha_{\mathbb F_2}\) requires checking that convention. The record also does not specify whether graphs must be simple or may have loops and parallel edges. These choices remain unresolved in this statement-only draft.
\subsection{Short English statement}
Can the least total length of cycles whose attached disks make a finite connected graph simply connected exceed, by an arbitrarily large factor, the least total length needed to eliminate first homology over \(\mathbb F_2\)? Can this happen with a single uniform bound on vertex degrees? Trees are excluded from the ratio, and the precise graph and attaching-cycle conventions require confirmation.
\subsection{Sources}
\begin{itemize}
\item Supplied local UnsolvedMath registry record, permanent ID \(30005353\), title \emph{Homotopical versus Homological Cycle-Filling Complexity}, source problem code \texttt{catalogue entry}, in the Oberwolfach Reports Open Problems collection; proposed year recorded as \(2023\). This draft uses the supplied complete record without independently inspecting the registry files.
\item Supplied report citation: \emph{Combinatorics} (2023), Oberwolfach Reports, DOI \texttt{10.4171/owr/2023/1}, pp.\ 5--89. Supplied source URL and literature-source link, labeled ``Oberwolfach Reports 2023/1'': \url{https://doi.org/10.4171/owr/2023/1}. The supplied background describes Oberwolfach Reports as documenting mathematical workshops and their open-problem discussions.
\item Supplied dated literature curation, checked \texttt{2026-08-22}: the record labels the problem open and reports that known constructions give additive separation between the two quantities, while no result establishes an unbounded ratio. This is supplied dated curation, not an independently verified literature assessment. The background explicitly requests independent verification of time-sensitive status and cited claims.
\item The supplied statement status is \texttt{corrected\_verified}; its verification field reports checking the clean statement against the official report and preserving the report's hypotheses, notation, and target. These are supplied verification claims; this draft does not independently verify them.
\item Attribution remains uncertain. The proposer field is empty. The raw transcription places the name Jacob Fox after the ratio question, followed by a statement about work of Conlon, Fox, Pham, and Yepremyan and a separate ``Question 11'' on monochromatic linear subspaces, including a reference to a conjecture of Noga Alon. The transcription does not reliably establish which attribution belongs to the present question.
\end{itemize}
\end{document}
